\documentclass[12pt]{article}
\usepackage{fullpage}
\usepackage{amsmath,amssymb,amsthm}

\newtheorem{lemma}{Lemma}
\newtheorem{proposition}{Proposition}
\newcommand{\R}{\mathbb R}
\newcommand{\E}{\mathbb E}
\newcommand{\1}{\mathbf 1}
\newcommand{\Var}{\operatorname{Var}}
\newcommand{\TV}{\operatorname{TV}}
\newcommand{\Law}{\operatorname{Law}}
\newcommand{\calF}{\mathcal F}
\newcommand{\Ent}{\operatorname{Ent}}
\newcommand{\Lip}{\operatorname{Lip}}

\begin{document}

\begin{center}
{\bf Uniqueness of the invariant measure}
\end{center}

\section*{Problem statement and interpretation}

Put \(E=C_0([0,1])\), \(H=L^2(0,1)\), and let \(S_t\) be the Dirichlet heat
semigroup generated by \(\frac12\Delta\).  I use the Dirichlet version of the
Athreya--Butkovsky--L\^e--Mytnik construction in the following precise form.
For \(b_n=G_{1/n}\), the regularized equation has the mild solution
\[
  u_t^n=S_t x+V_t+K_t^n,\qquad
  K_t^n=\int_0^tS_{t-r}b_n(u_r^n)\,dr,\qquad
  V_t=\int_0^tS_{t-r}\,dW_r ,
\]
and, for every \(t>0\) and every \(\Phi\in C_b(E)\),
\[
        P_t^n\Phi(x):=\E\Phi(u_t^n)\longrightarrow P_t\Phi(x).
        \tag{1}
\]
This is exactly the weak convergence assumption of the question, written on
the Polish state space \(E\).

\begin{lemma}[Uniform estimate on the regularized drift]\label{lem:Hestimate}
For each \(T<\infty\) there is \(C_T<\infty\), independent of \(n\) and of
\(x\in E\), such that
\[
 \|K_t^n-S_{t-s}K_s^n\|_{L^2(\Omega;H)}
      \le C_T |t-s|^{3/4},\qquad 0\le s<t\le T .             \tag{2}
\]
\end{lemma}

\begin{proof}
Let \(p_t\) be the Dirichlet kernel and \(d(y)=y\wedge(1-y)\).  For
\(0<a\le T\),
\[
p_a(z,y)\le Ca^{-1/2}e^{-|z-y|^2/(Ca)},\qquad
p_a(y,y)\asymp a^{-1/2}\bigl(1\wedge d(y)^2/a\bigr).
\]
The diagonal estimate is the usual one-dimensional boundary estimate; it
follows, for example, from the reflection formula for the interval, compared
with the half-line identity
\((2\pi a)^{-1/2}(1-e^{-2d(y)^2/a})\).  Consequently, for
\(\rho_\tau(y)=\int_0^\tau p_{2a}(y,y)\,da\),
\[
        \rho_\tau(y)\asymp \sqrt{\tau}\wedge d(y),             \tag{3}
\]
with constants depending only on \(T\).  Hence
\[
 \int_0^1\rho_\tau(y)^{-1/2}dy\le C_T\tau^{-1/4}.              \tag{4}
\]
Also
\[
 \sup_z\int_0^1p_a(z,y)\rho_\tau(y)^{-1/2}dy
 \le C_T(a^{-1/4}+\tau^{-1/4}).                              \tag{5}
\]
Indeed, \(\rho_\tau^{-1/2}\le C(\tau^{-1/4}+d^{-1/2})\).  The
\(\tau^{-1/4}\) part integrates to \(\tau^{-1/4}\).  For the other part,
\(d(y)^{-1/2}\le y^{-1/2}+(1-y)^{-1/2}\), and after the change
\(y=\sqrt a\,r\) (or \(1-y=\sqrt a\,r\)) the remaining bound is
\(Ca^{-1/4}\sup_{\alpha\ge0}\int_0^\infty e^{-(r-\alpha)^2/C}r^{-1/2}dr\),
which is finite.  Finally,
\[
        \|\rho_\tau^{-1}S_\tau f\|_2\le C_T\tau^{-1/2}\|f\|_2, \tag{6}
\]
because \(\rho_\tau^{-1}\le C(\tau^{-1/2}+d^{-1})\), and Hardy's inequality
\(\|g/d\|_2\le C\|g'\|_2\), applied to \(g=S_\tau f\), is combined with
\(\|\partial_xS_\tau f\|_2\le C\tau^{-1/2}\|f\|_2\).

Write \(D^n_{s,t}=K_t^n-S_{t-s}K_s^n\).  Freeze the equation at time \(s\):
\[
 A^n_{s,t}=\int_s^tS_{t-r}
 b_n\!\left(S_{r-s}u_s^n+V_r-S_{r-s}V_s\right)\,dr .          \tag{7}
\]
For fixed \(n\), if \(\Pi\) is a partition of \([s,t]\), then
\[
 \sum_{[a,b]\in\Pi}S_{t-b}A^n_{a,b}\to D^n_{s,t}\quad\hbox{in }L^2(H).
 \tag{8}
\]
This follows by comparing \(A^n_{a,b}\) with
\[
        \int_a^bS_{b-r}b_n(u_r^n)\,dr .
\]
The difference is at most
\(\Lip(b_n)\|b_n\|_\infty(b-a)^2\) in \(L^2(H)\), so the sum of the errors is
\(O((t-s)|\Pi|)\).

We estimate \(A^n\), uniformly in \(n\).  Given \(\calF_s\),
\(Z_r(y)=S_{r-s}u_s^n(y)+V_r(y)-S_{r-s}V_s(y)\) is Gaussian with variance
\(\rho_{r-s}(y)\).  Since \(b_n\) is a probability density, convolution with a
Gaussian of variance \(\sigma^2\) is bounded by \(C\sigma^{-1}\).  Conditioning
at time \(r\) for the frozen Gaussian process gives, for \(s<r<q\),
\[
 \E_s\,b_n(Z_r(y))b_n(Z_q(z))
 \le C\rho_{r-s}(y)^{-1/2}\rho_{q-r}(z)^{-1/2}.              \tag{9}
\]
Using \(\int p_{t-r}(x,y)p_{t-q}(x,z)dx=p_{2t-r-q}(y,z)\), (4), and (5),
\[
 \E_s\|A^n_{s,t}\|_2^2
 \le C\!\int_s^t\!\!\int_r^t
      \bigl((2t-r-q)^{-1/4}+(r-s)^{-1/4}\bigr)(q-r)^{-1/4}\,dq\,dr
 \le C|t-s|^{3/2}.                                          \tag{10}
\]
Thus \(\|A^n_{s,t}\|_{L^2H}\le C|t-s|^{3/4}\).

For \(s<u<t\), set
\(\delta A^n_{s,u,t}=A^n_{s,t}-S_{t-u}A^n_{s,u}-A^n_{u,t}\).  In the part
after \(u\), the two frozen arguments differ by \(h=S_{r-u}D^n_{s,u}\).
Since \(\|(b_n*\varphi_\sigma)'\|_\infty\le C\sigma^{-2}\) uniformly in \(n\),
\[
 \left|\E_u\{b_n(Z_r(y))-b_n(Z_r(y)+h(y))\}\right|
      \le C\rho_{r-u}(y)^{-1}|h(y)|.
\]
By (6),
\[
        \|\E_u\delta A^n_{s,u,t}\|_2
        \le C(t-u)^{1/2}\|D^n_{s,u}\|_2 .                    \tag{11}
\]

We shall use the following elementary Hilbert-valued sewing estimate.  Let
\(A_{s,t}\) be \(H\)-valued, \(\calF_t\)-measurable, and let
\(S_t\) be a contraction semigroup.  For a dyadic interval \([s,t]\), denote
by \(J_m=\sum_iS_{t-t_{i+1}}A_{t_i,t_{i+1}}\) the sum over the \(2^m\) equal
subintervals.  If for every dyadic child \([a,b]\) with midpoint \(u\)
\[
\|\E_a\delta A_{a,u,b}\|_{L^2H}\le \Gamma_1(b-a)^{1+\varepsilon_1},\qquad
\|\delta A_{a,u,b}-\E_a\delta A_{a,u,b}\|_{L^2H}
        \le \Gamma_2(b-a)^{1/2+\varepsilon_2},
\]
then \(J_m\) converges in \(L^2(H)\) and its limit \({\cal A}_{s,t}\) satisfies
\[
 \|{\cal A}_{s,t}-A_{s,t}\|_{L^2H}
 \le C\bigl(\Gamma_1|t-s|^{1+\varepsilon_1}
       +\Gamma_2|t-s|^{1/2+\varepsilon_2}\bigr).             \tag{12}
\]
Indeed, \(J_{m+1}-J_m\) is the sum of
\(-S_{t-b}\delta A_{a,u,b}\).  The conditional-mean parts have norm bounded by
\(C\Gamma_1|t-s|^{1+\varepsilon_1}2^{-m\varepsilon_1}\).  The centered parts
are orthogonal martingale differences: if two dyadic intervals are ordered,
the earlier summand is measurable before the conditioning time of the later
one.  Hence their square norm is bounded by
\(C\Gamma_2^2|t-s|^{1+2\varepsilon_2}2^{-2m\varepsilon_2}\).  The geometric
series gives (12).

For fixed \(n\), \(Y_n(h)=\sup_{0<t-s\le h}\|D^n_{s,t}\|_{L^2H}/|t-s|^{3/4}\)
is finite because \(b_n\) is bounded.  On intervals of length at most \(h\),
(11) gives
\[
 \|\E_a\delta A^n_{a,u,b}\|_{L^2H}\le C Y_n(h)(b-a)^{5/4},
\]
while (10) gives the centered bound \(C(b-a)^{3/4}\).  Combining (8), (10) and
(12),
\[
 \|D^n_{s,t}\|_{L^2H}\le C|t-s|^{3/4}
        +C Y_n(h)|t-s|^{5/4}.
\]
Thus \(Y_n(h)\le C+Ch^{1/2}Y_n(h)\).  Choosing \(h\) small absorbs the second
term.  Finally use
\(D^n_{s,t}=S_{t-r}D^n_{s,r}+D^n_{r,t}\) and split a general interval into
pieces of length at most \(h\).  This proves (2).
\end{proof}

\section*{The reversible reference measure and its gap}

Let \(\gamma\) be Brownian bridge measure on \(E\), i.e. the centered Gaussian
measure with covariance \((-\Delta_D)^{-1}\).  Define
\[
  \Lambda(v)=\int_0^1\1_{\{v(r)>0\}}\,dr,\qquad
  \pi(dv)=Z^{-1}e^{2\Lambda(v)}\,\gamma(dv).                 \tag{13}
\]
For \(B_n'=b_n\), \(B_n(a)=\int_{-\infty}^ab_n(r)dr\), put
\[
   \pi_n(dv)=Z_n^{-1}\exp\left(2\int_0^1B_n(v(r))\,dr\right)\gamma(dv).
\]
Since \(0\le B_n\le1\), \(B_n(a)\to\1_{\{a>0\}}\) for \(a\ne0\), and a
Brownian bridge spends zero Lebesgue time at a fixed level,
\[
        \|\pi_n-\pi\|_{\TV}\to0 .                            \tag{14}
\]
Moreover all densities \(d\pi_n/d\gamma\) and \(d\pi/d\gamma\) lie between
two deterministic positive constants.

\begin{lemma}\label{lem:reggap}
There are constants \(C,c>0\), independent of \(n\), such that for every
bounded continuous \(f,g\) and every \(t\ge0\),
\[
 \int P_t^n f\,d\pi_n=\int f\,d\pi_n,\qquad
 \int gP_t^n f\,d\pi_n=\int fP_t^n g\,d\pi_n,                \tag{15}
\]
and
\[
        \Var_{\pi_n}(P_t^n f)\le e^{-ct}\Var_{\pi_n}(f).      \tag{16}
\]
\end{lemma}

\begin{proof}
Let \(\Pi_m\) be projection onto the first \(m\) Dirichlet modes and set
\[
        U_{n,m}(v)=\int_0^1 B_n(\Pi_m v(r))\,dr .
\]
Let \(u^{n,m}\) solve
\[
u_t^{n,m}=S_t x+V_t+\int_0^tS_{t-r}\Pi_m b_n(\Pi_m u_r^{n,m})\,dr
\]
and denote its semigroup by \(P_t^{n,m}\).  The drift is the \(H\)-gradient of
\(U_{n,m}\).  Thus the first \(m\) Fourier coordinates form a smooth
finite-dimensional Langevin diffusion with invariant density
\(Z_{n,m}^{-1}e^{2U_{n,m}}\) relative to their Gaussian law, and the remaining
coordinates are an independent OU tail.  Finite-dimensional integration by
parts, together with the standard OU integration by parts in the tail, gives
invariance and symmetry of \(P_t^{n,m}\) with respect to
\(\pi_{n,m}=Z_{n,m}^{-1}e^{2U_{n,m}}\gamma\).

The measures \(\pi_{n,m}\) converge to \(\pi_n\) in total variation: indeed
\(\Pi_m v\to v\) in \(H\), \(\gamma\)-a.s., and \(B_n\) is bounded and
Lipschitz.  Also, for fixed \(n\), \(u_t^{n,m}\to u_t^n\) in probability in
\(E\), for every \(t>0\).  To see this, first work in \(H\).  The Lipschitz
property gives, after splitting
\(\Pi_m b_n(\Pi_m z)-b_n(y)\) into the three errors
\(\Pi_m[b_n(\Pi_m z)-b_n(\Pi_m y)]\),
\(\Pi_m[b_n(\Pi_m y)-b_n(y)]\), and \((\Pi_m-I)b_n(y)\), a Gronwall estimate;
the last two errors vanish by dominated convergence in
\(L^1([0,T]\times\Omega;H)\).  For convergence in \(E\) at a fixed positive
time, split the drift convolution over \([0,t-\varepsilon]\) and
\([t-\varepsilon,t]\); the first part is smoothed by \(S_{t-r}\) into
\(H^\alpha_0\subset E\), \(\alpha>1/2\), and the second is uniformly small
because \(b_n\) is bounded.

Passing \(m\to\infty\) in the reversible identities therefore proves (15) for
bounded continuous \(f,g\).  The Gaussian Poincare inequality, with the
\(H\)-gradient, and the two-sided bound on \(e^{2U_{n,m}}/Z_{n,m}\), give by
the Holley--Stroock perturbation argument
\[
      \Var_{\pi_{n,m}}(F)\le C\,{\cal E}_{n,m}(F,F)
\]
with \(C\) independent of \(m,n\).  Since \(P^{n,m}\) is symmetric, this
Poincare inequality gives
\[
 \Var_{\pi_{n,m}}(P_t^{n,m}f)\le e^{-ct}\Var_{\pi_{n,m}}(f).
\]
Letting \(m\to\infty\), using the above convergence and the uniformly bounded
densities, yields (16).
\end{proof}

\begin{proposition}\label{prop:gap}
The measure \(\pi\) is \(P_t\)-invariant, \(P_t\) is symmetric on
\(L^2(\pi)\), and
\[
       \Var_\pi(P_t f)\le e^{-ct}\Var_\pi(f),
       \qquad f\in L^2(\pi),\ t\ge0 .                       \tag{17}
\]
\end{proposition}

\begin{proof}
For \(f,g\in C_b(E)\), (1), (14), (15), and the uniform boundedness of the
densities imply
\[
 \int P_tf\,d\pi=\lim_n\int P_t^nf\,d\pi_n=\int f\,d\pi,\qquad
 \int gP_tf\,d\pi=\int fP_tg\,d\pi .
\]
On the Polish space \(E\), these identities identify respectively the finite
measure \(\int\pi(dx)P_t(x,\cdot)\) with \(\pi\), and the measure
\(\pi(dx)P_t(x,dy)\) with its transpose.  Hence invariance and detailed balance
hold for all bounded Borel functions.  Likewise (16) passes to the limit for
bounded continuous \(f\).  Since \(C_b(E)\) is dense in \(L^2(\pi)\), and an
invariant Markov operator is an \(L^2\)-contraction, symmetry and (17) extend
to all \(L^2(\pi)\).
\end{proof}

\section*{Fixed-time absolute continuity}

\begin{lemma}\label{lem:control}
For every deterministic \(x\in E\) and every \(t>0\),
\[
        P_t(x,\cdot)\ll \gamma ,
\]
and hence \(P_t(x,\cdot)\ll\pi\).
\end{lemma}

\begin{proof}
Let \(\nu_t^x=\Law(S_tx+V_t)\).  We prove first that
\(\Ent(P_t^n(x,\cdot)\mid\nu_t^x)\) is bounded uniformly in \(n\).  Set
\(r_k=t(1-2^{-k})\), \(\Delta_k=r_{k+1}-r_k\), and
\(D_k^n=K^n_{r_{k+1}}-S_{\Delta_k}K^n_{r_k}\).  Define
\[
 F_s^n=\sum_{k\ge0}\frac{\1_{(r_{k+1},r_{k+2}]}(s)}{\Delta_{k+1}}\,
        S_{s-r_{k+1}}D_k^n .                                \tag{18}
\]
This process is predictable, because \(D_k^n\) is \(\calF_{r_{k+1}}\)-measurable.
By Lemma \ref{lem:Hestimate},
\[
 \E\int_0^t\|F_s^n\|_2^2ds
 \le C_t\sum_{k\ge0}\frac{\Delta_k^{3/2}}{\Delta_{k+1}}<\infty, \tag{19}
\]
uniformly in \(n\), and its terminal effect telescopes:
\[
 \int_0^tS_{t-s}F_s^n\,ds
 =\sum_{k\ge0}S_{t-r_{k+1}}D_k^n=K_t^n                       \tag{20}
\]
in \(H\).  Moreover, for every \(F\in L^2(0,t;H)\),
\(\int_0^tS_{t-s}F_sds\in H^\alpha_0(0,1)\subset E\) for
\(\alpha\in(1/2,1)\), by the Dirichlet-basis estimate
\[
 \sum_j\lambda_j^\alpha\Big|\int_0^te^{-\lambda_j(t-s)/2}F_j(s)ds\Big|^2
 \le \sum_j\lambda_j^{\alpha-1}\int_0^t|F_j(s)|^2ds .
\]

To apply Girsanov rigorously, realize the cylindrical Brownian motion as a
continuous process in a Hilbert space \(U\) for which the embedding
\(H\hookrightarrow U\) is Hilbert--Schmidt, for instance \(U=H^{-\beta}\) with
\(\beta>1/2\).  The standard Girsanov theorem on
\(C([0,t];U)\) says that an adapted \(H\)-valued drift \(F\) with
\(\int_0^t\|F_s\|_H^2ds<\infty\) after localization changes the noise law with
relative entropy at most one half of its energy.  Let
\(\tau_N=\inf\{s:\int_0^s\|F_r^n\|_2^2dr>N\}\) and
\(F^{n,N}=F^n\1_{[0,\tau_N]}\).  For
\[
 X_N=S_tx+V_t+\int_0^tS_{t-s}F_s^{n,N}ds ,
\]
Novikov's condition holds, and the data-processing inequality gives
\[
 \Ent(\Law(X_N)\mid\nu_t^x)
 \le \frac12\,\E\int_0^t\|F_s^{n,N}\|_2^2ds
 \le \frac12\,\E\int_0^t\|F_s^n\|_2^2ds .                   \tag{21}
\]
On \(\{\tau_N>t\}\), (20) gives \(X_N=u_t^n\), and
\(\mathbb P(\tau_N\le t)\le N^{-1}\E\int_0^t\|F_s^n\|_2^2ds\to0\).  Hence
\(X_N\to u_t^n\) in probability in \(E\).  Lower semicontinuity of relative
entropy on \(E\), with (19)--(21), yields
\[
        \sup_n\Ent(P_t^n(x,\cdot)\mid\nu_t^x)<\infty .       \tag{22}
\]
Using the weak convergence (1) and lower semicontinuity once more,
\(\Ent(P_t(x,\cdot)\mid\nu_t^x)<\infty\); hence
\(P_t(x,\cdot)\ll\nu_t^x\).

Finally compare \(\nu_t^x\) with \(\gamma\).  In the basis
\(e_j(r)=\sqrt2\sin(j\pi r)\), the covariance of \(\nu_t^x\) has eigenvalues
\[
        q_j(t)=\frac{1-e^{-j^2\pi^2t}}{j^2\pi^2},
\]
whereas \(\gamma\) has eigenvalues \((j^2\pi^2)^{-1}\).  The ratios are bounded
away from zero and differ from \(1\) by a square-summable sequence; moreover
\(S_tx\in H^1_0(0,1)\), the Cameron--Martin space of \(\gamma\).  By the
Feldman--Hajek theorem for Gaussian measures on \(E\), \(\nu_t^x\sim\gamma\).
\end{proof}

\section*{Conclusion}

Let \(\eta\) be an invariant probability measure.  If \(\pi(A)=0\), then
Lemma \ref{lem:control} gives \(P_t(y,A)=0\) for every \(y\in E\), and hence
\[
        \eta(A)=\int_E P_t(y,A)\,\eta(dy)=0.
\]
Thus \(\eta\ll\pi\); write \(d\eta=h\,d\pi\).  Detailed balance gives the
\(L^1\)-duality
\[
        \int f\,P_t g\,d\pi=\int P_t f\,g\,d\pi
\]
for bounded Borel \(f\) and \(g\in L^1(\pi)\), by truncating \(g\) and using
the bounded Borel identity.  Therefore
\[
 \int f h\,d\pi=\int P_tf\,d\eta
      =\int P_tf\,h\,d\pi=\int f\,P_t h\,d\pi ,
\]
so \(P_t h=h\) in \(L^1(\pi)\).  For \(a>0\), let \(h_a=h\wedge a\).  Since
\(r\mapsto r\wedge a\) is concave,
\[
        P_t h_a\le (P_t h)\wedge a=h_a .
\]
The two sides have the same \(\pi\)-integral; therefore \(P_t h_a=h_a\).
Now \(h_a\in L^2(\pi)\), and the spectral gap (17) implies
\(\Var_\pi(h_a)=0\).  Thus \(h_a\) is constant for every \(a\).  Letting
\(a\uparrow\infty\) and using \(\int h\,d\pi=1\), we obtain \(h=1\)
\(\pi\)-a.s.  Hence \(\eta=\pi\).

Therefore \(P_t\) has at most one invariant probability measure on
\(C_0([0,1])\).  In fact, if an invariant probability exists, it is the Gibbs
measure \(\pi\) in (13).

\begin{thebibliography}{9}

\bibitem{ABLM}
S. Athreya, O. Butkovsky, K. L\^e and L. Mytnik,
\emph{Well-posedness of stochastic heat equation with distributional drift
and skew stochastic heat equation}, Comm. Pure Appl. Math. 77 (2024),
2708--2777.

\bibitem{LeBanach}
K. L\^e,
\emph{Stochastic sewing in Banach spaces}, Electron. J. Probab. 28 (2023),
paper no. 26, 1--22.

\bibitem{LeSSL}
K. L\^e,
\emph{A stochastic sewing lemma and applications}, Electron. J. Probab. 25
(2020), paper no. 38, 1--55.

\bibitem{Bogachev}
V. I. Bogachev,
\emph{Gaussian Measures}, American Mathematical Society, 1998.

\bibitem{DPZ}
G. Da Prato and J. Zabczyk,
\emph{Ergodicity for Infinite Dimensional Systems}, Cambridge University
Press, 1996.

\bibitem{MR}
Z.-M. Ma and M. R\"ockner,
\emph{Introduction to the Theory of (Non-Symmetric) Dirichlet Forms},
Springer, 1992.

\bibitem{BZ}
S. K. Bounebache and L. Zambotti,
\emph{A skew stochastic heat equation}, J. Theoret. Probab. 27 (2014),
168--201.

\end{thebibliography}

\end{document}
